\input pdfToolbox
\input preamble
\input pdfmsym

\pdfmsymsetscalefactor{10}

\setlayout{horizontal margin=2cm, vertical margin=2cm}
\parindent=0pt
\parskip=3pt plus 2pt minus 2pt

\input preamble

\footline={}

%%%%%%%%%%%%%%%%%%%%%%%%%%%%%%%%%%%%%%%%%%%%%%%%%%%%%%%%%%%%%%%%

\headline={\pageborder{rgb{1 .8 .8}}{rgb{.6 0 0}}{5}}

\color rgb{1 .5 .1}

{\def\boxshadowcolor{rgb{1 .7 .7}}
\bppbox{rgb{1 .8 .8}}{rgb{.6 0 0}}{rgb{.2 0 0}}

    \centerline{\setfontandscale{bf}{20pt}Category Theory}
    \smallskip
    \centerline{\setfont{it}Summary by Ari Feiglin \setfont{rm}({\tt ari.feiglin@gmail.com})}

\eppbox

\bigskip

\bppbox{rgb{1 .8 .8}}{rgb{.6 0 0}}{rgb{.2 0 0}}
    \section*{Contents}
    
    \tableofcontents
\eppbox

}

\vfill\break

\color{black}

\pageno=1
\newif\ifpageodd
\pageoddtrue
\headline={%
    \hbox to \hsize{\color{black}%
        \ifpageodd\hfil{\it\currsubsection\quad\bf\folio}\global\pageoddfalse%
        \else{\bf\folio\quad\it\currsubsection}\hfil\global\pageoddtrue\fi%
    }%
}

I begin by noting that these notes are for those already familiar with the very basic definitions of category theory (those of categories, functors, and natural transformations).

\input equivalences
\input limits
\input yoneda
\input adjunctions
\input abelian-categories

\bye

